\documentclass[10pt,a4paper]{amsart}
\usepackage[margin=19mm]{geometry}
\usepackage{amsmath,amssymb,mathtools}
\usepackage[scr=rsfs,cal=euler]{mathalfa}
\usepackage{xfrac}
\usepackage[unicode,hidelinks,bookmarks=false]{hyperref}
\newcommand{\ie}{i.e.\ }
\newcommand{\cf}{cf.\ }
\newcommand{\resp}{respectively\ }
\newcommand{\Subsection}{\subsection}
\providecommand{\nobreakdash}{\nobreak\hbox{-}}
\setlength{\emergencystretch}{2em}
\multlinegap0pt
\usepackage{amssymb}
\usepackage{algorithm}\usepackage{algpseudocode}
\usepackage{xcolor}
\usepackage{tcolorbox}
\usepackage{tikz}
\newtheorem{prop}{Proposition}
\title{Sequential learning theory for Markov genealogy processes}
\author{David J Pascall}
\date{Source: arXiv:2603.09033v2 (CC BY 4.0)}
\begin{document}
\maketitle

\noindent\emph{Source and licence.} Sequential learning theory for Markov genealogy processes. Version arXiv:2603.09033v2 (CC BY 4.0), distributed by arXiv under the Creative Commons Attribution 4.0 International licence, http://creativecommons.org/licenses/by/4.0/, as recorded in the arXiv licence metadata for this exact version and shown on the abstract page https://arxiv.org/abs/2603.09033v2. Full licence notice: \texttt{licenses/arxiv-2603.09033-CC-BY-4.0.txt}.\par\smallskip

\noindent\textit{Editorial scope.} Counted unit: the article's prop metricstandardlearning (source0.tex lines 302--305). The article's complete original proof of it is delivered with it (source0.tex lines 306--312, one \texttt{proof} environment of 7 lines). No mathematics is written, repaired or re-derived: the statement and the proof are the article's own lines, in the article's own notation, and no retained line is substituted or re-flowed.  No further source-local context is needed: every cross-reference of every retained line resolves inside the two delivered spans.  Nothing was omitted from any retained span, so the package carries no omission record.  No retained line contains a citation, so the wrapper carries no bibliography and no bibliography entry.  No retained line contains a figure environment, an image-inclusion command or drawing code, so no image is delivered and the package declares no asset dependency.  Editorial formatting only, in addition: an AMS \texttt{amsart} preamble that restates the article's own declarations and macros used by the retained lines, namely the environments \texttt{prop} on separate counters, together with the standard packages those lines need.  The retained environments print in this document as Proposition 1 in order of appearance, whereas the article prints the same items with the numbering of its own full document.  The complete native source of the article as distributed in the arXiv e-print for this exact version is bundled unchanged as \texttt{sources/196041/source0.tex}.
\begin{prop}[Standard learning for fixed metric-space-valued estimands]
\label{metricstandardlearning}
Let $(S,d)$ be a separable metric space and let $K$ be a permutation-invariant $S$-valued estimand satisfying $\mathbb{E}_{\mathbb{P}^*}[d^2(K,s_0)]<\infty$ for some $s_0 \in S$. Then the conditional Fr\'{e}chet variance of $K$ given $\mathcal{F}_n$, $\inf_{s \in S} \mathbb{E}_{\mathbb{P}^*}[d^2(K,s)|\mathcal{F}_n]$, is a supermartingale. That is, for all $n < f(\mathcal{G})$, $\mathbb{E}_{\mathbb{P}^*}[\inf_{s \in S} \mathbb{E}_{\mathbb{P}^*}[d^2(K,s)|\mathcal{F}_{n+1}]|\mathcal{F}_n] \leq \inf_{s \in S} \mathbb{E}_{\mathbb{P}^*}[d^2(K,s)|\mathcal{F}_n]$ a.s.
\end{prop}

\begin{proof}
\textit{Well-posedness.} For each $s \in S$, the map $x \mapsto d^2(x,s)$ is continuous, hence $d^2(K,s)$ is $\mathcal{F}$-measurable. Moreover, for any $s \in S$, $d^2(K,s) \leq 2d^2(K,s_0) + 2d^2(s_0,s)$, so $d^2(K,s) \in L^1$ by the basepointed second-moment assumption. Thus $\mathbb{E}_{\mathbb{P}^*}[d^2(K,s)|\mathcal{F}_n]$ is a finite $\mathcal{F}_n$-measurable random variable for each $s$.

For $s, s' \in S$, $|d^2(K,s) - d^2(K,s')| \leq d(s,s')(2d(K,s_0) + d(s_0,s) + d(s_0,s'))$. Since $d(K,s_0) \in L^1$ by Jensen and the basepointing assumption, taking conditional expectations yields $|\mathbb{E}_{\mathbb{P}^*}[d^2(K,s)|\mathcal{F}_n] - \mathbb{E}_{\mathbb{P}^*}[d^2(K,s')|\mathcal{F}_n]| \leq d(s,s')(2\mathbb{E}_{\mathbb{P}^*}[d(K,s_0)|\mathcal{F}_n] + d(s_0,s) + d(s_0,s')) \quad \text{a.s.}$ for each fixed pair $(s,s')$. Fix a countable dense subset $\{s_k\} \subset S$. Restricting to pairs from $\{s_k\}$, the bound holds simultaneously for all pairs outside a single null set; the resulting map extends uniquely to a continuous version of $s \mapsto \mathbb{E}_{\mathbb{P}^*}[d^2(K,s)|\mathcal{F}_n](\omega)$ on $S$. Working with this version, $\inf_{s \in S} \mathbb{E}_{\mathbb{P}^*}[d^2(K,s)|\mathcal{F}_n] = \inf_k \mathbb{E}_{\mathbb{P}^*}[d^2(K,s_k)|\mathcal{F}_n]$ a.s., so the conditional Fr\'{e}chet variance is $\mathcal{F}_n$-measurable as a countable infimum of $\mathcal{F}_n$-measurable random variables. Finally, $0 \leq \inf_{s \in S} \mathbb{E}_{\mathbb{P}^*}[d^2(K,s)|\mathcal{F}_n] \leq \mathbb{E}_{\mathbb{P}^*}[d^2(K,s_0)|\mathcal{F}_n] < \infty$ a.s., implying $\inf_{s \in S} \mathbb{E}_{\mathbb{P}^*}[d^2(K,s)|\mathcal{F}_n] \in L^1$.

\textit{Inequality.} For any fixed $s \in S$, the tower property gives $\mathbb{E}_{\mathbb{P}^*}[d^2(K,s)|\mathcal{F}_n] = \mathbb{E}_{\mathbb{P}^*}[\mathbb{E}_{\mathbb{P}^*}[d^2(K,s)|\mathcal{F}_{n+1}]|\mathcal{F}_n] \geq \mathbb{E}_{\mathbb{P}^*}[\inf_{s' \in S}\mathbb{E}_{\mathbb{P}^*}[d^2(K,s')|\mathcal{F}_{n+1}]|\mathcal{F}_n]$, where the inequality is pointwise: for each $\omega$, the conditional expectation at $s$ is at least the infimum over $s'$. As the right-hand side is independent of $s$, taking the infimum over $s$ on the left gives the result.
\end{proof}

\end{document}
